@@M0@@
\uline{is
exact}.

If @@M1@@ is a coherent @@M2@@-Module, @@M3@@ is an @@M4@@-Module whose support is contained in @@M5@@; consequently its restriction to @@M6@@ is zero, and @@M7@@, whence, since the elements of the form @@M8@@ generate @@M9@@, @@M10@@. This gives a homomorphism
@@M11@@
We shall define an inverse homomorphism to @@M12@@. Let @@M13@@ be a coherent @@M14@@-Module. Since @@M15@@ is noetherian, there exists a coherent @@M16@@-Module @@M17@@ whose restriction to @@M18@@ is isomorphic to @@M19@@. If @@M20@@ is a second coherent @@M21@@-Module whose restriction to @@M22@@ is isomorphic to @@M23@@, then we have in @@M24@@:
@@M25@@
To see this, regard @@M26@@ as a subsheaf of @@M27@@ by means of the diagonal map; @@M28@@ then extends to a subsheaf @@M29@@ of @@M30@@, and we are reduced to showing that
@@M31@@
and the analogous assertion for @@M32@@. Let then @@M33@@ and @@M34@@ be the kernel and cokernel of the composite homomorphism:
@@M35@@
this homomorphism reduces to the identity of @@M36@@ over @@M37@@, and therefore @@M38@@ and @@M39@@ have support contained in @@M40@@. If @@M41@@ is the ideal of @@M42@@ in @@M43@@, it follows that there exist integers @@M44@@ and @@M45@@ such that @@M46@@, @@M47@@ (EGA I 9.3.4). One may therefore express @@M48@@ and @@M49@@ as successive extensions of Modules
